% Generated by code/make_cue_cards.py from docs/cue_cards.md -- edit the Markdown, not this file.
\documentclass[10pt]{article}
\usepackage[a4paper,landscape,margin=0.7cm]{geometry}
\usepackage[T1]{fontenc}
\usepackage[utf8]{inputenc}
\usepackage[sfdefault]{FiraSans}
\usepackage{xcolor,newunicodechar,pifont,enumitem}
\usepackage[most]{tcolorbox}
\tcbuselibrary{fitting}
\definecolor{navy}{HTML}{1F3B73}\definecolor{nlred}{HTML}{C0392B}\definecolor{grey}{HTML}{7F8C8D}\definecolor{gold}{HTML}{E0A100}
\newunicodechar{−}{\textminus}\newunicodechar{→}{\ensuremath{\rightarrow}}\newunicodechar{≈}{\ensuremath{\approx}}
\newunicodechar{≤}{\ensuremath{\leq}}\newunicodechar{↑}{\ensuremath{\uparrow}}\newunicodechar{€}{\texteuro}
\newunicodechar{✔}{\textcolor{green!50!black}{\ding{51}}}\newunicodechar{★}{\textcolor{gold}{\ding{72}}}
\newunicodechar{…}{\ldots}
\setlist[itemize]{leftmargin=1.1em,itemsep=1pt,topsep=2pt,parsep=0pt,label=\textcolor{#1}{\textbullet}}
\pagestyle{empty}\setlength{\parindent}{0pt}
\newtcolorbox{card}[2]{enhanced,nobeforeafter,box align=top,width=0.492\linewidth,height=0.475\textheight,fit,fit fontsize macros,fit basedim=17pt,fit algorithm=fontsize,
  colframe=#1,colback=white,boxrule=0.6pt,arc=2pt,left=6pt,right=6pt,top=4pt,bottom=4pt,
  borderline={0.3pt}{-3.5pt}{black!35,dashed},
  title={#2},fonttitle=\bfseries\large,coltitle=white,colbacktitle=#1,toptitle=2pt,bottomtitle=2pt}
\begin{document}

\begin{card}{navy}{\makebox[0pt][l]{0 \quad Title}\hfill{\normalsize\mdseries Zheng · 0:15 · from 0:00}}
{\large\itshape "We are Zheng Huang and Diego Gutiérrez."}\par\smallskip
\begin{itemize}[label=\textcolor{navy}{\textbullet}]\item Case: the Netherlands, CCyB 0 → 2\%\item How and why DNB did it\end{itemize}
\vfill\textcolor{grey}{\rule{\linewidth}{0.3pt}}\par{\bfseries\color{navy}→}\ "Here is the whole talk on one slide."
\end{card}
\hfill
\begin{card}{navy}{\makebox[0pt][l]{1 \quad One-pager}\hfill{\normalsize\mdseries Zheng · 0:45 · from 0:15}}
{\large\itshape "The whole talk on one slide."}\par\smallskip
\begin{itemize}[label=\textcolor{navy}{\textbullet}]\item DNB raised twice, to 2\%; the rulebook said \textbf{zero}\item Verdict: \textbf{normalisation, not tightening}\item Cheap · largely offset · \textbf{no sign of contraction}\end{itemize}
\vfill\textcolor{grey}{\rule{\linewidth}{0.3pt}}\par{\bfseries\color{navy}→}\ "Let's start in March 2020."
\end{card}


\vspace{0.35cm}

\begin{card}{navy}{\makebox[0pt][l]{2 \quad March 2020}\hfill{\normalsize\mdseries Zheng · 1:30 · from 1:00}}
{\large\itshape "March 2020. COVID has just hit Europe."}\par\smallskip
\begin{itemize}[label=\textcolor{navy}{\textbullet}]\item Neighbours \textbf{release} buffers within days (DK, NO, SE, DE, FR, IE…)\item \textcolor{gold!80!black}{\textbf{DEF}} buffer = extra equity above the minimum\item NL at \textbf{0\%}: nothing to release, so DNB improvised\item Systemic buffers cut 3\% → 2.5/2/1.5\% (ING, Rabo, ABN)\item \emph{(pause)} \textbf{The promise:} rebuild it as a 2\% CCyB\end{itemize}
\smallskip{\color{navy}\textbf{\#}}\ €8bn freed · up to €200bn of lending\par
\vfill\textcolor{grey}{\rule{\linewidth}{0.3pt}}\par{\bfseries\color{navy}→}\ "Why does it matter \emph{which} buffer you hold?"
\end{card}
\hfill
\begin{card}{navy}{\makebox[0pt][l]{3 \quad Releasable buffer}\hfill{\normalsize\mdseries Zheng · 1:00 · from 2:30}}
{\large\itshape "Not all capital is usable in a crisis."}\par\smallskip
\begin{itemize}[label=\textcolor{navy}{\textbullet}]\item CCyB = the only layer that can be \textbf{switched off}\item \textcolor{gold!80!black}{\textbf{DEF}} RWA = assets weighted by risk\item Other buffers: dipping in triggers \textbf{MDA} limits (dividends, bonuses)\item COVID: banks near thresholds \textbf{cut lending} anyway\item Key idea: \emph{how much} the authority can hand back\end{itemize}
\smallskip{\color{navy}\textbf{\#}}\ 0–2.5\% of RWA · up = 12 months · release = immediate\par
\vfill\textcolor{grey}{\rule{\linewidth}{0.3pt}}\par{\bfseries\color{navy}→}\ "Fast-forward to 2022, and look at the timing."
\end{card}

\newpage

\begin{card}{navy}{\makebox[0pt][l]{4 \quad The question}\hfill{\normalsize\mdseries Zheng · 1:00 · from 3:30}}
{\large\itshape "DNB keeps its promise, but look at the timing."}\par\smallskip
\begin{itemize}[label=\textcolor{navy}{\textbullet}]\item War · energy · inflation · ECB hikes · housing turning\item Basel indicator says \textbf{zero}\item \textcolor{gold!80!black}{\textbf{DEF}} procyclical = amplifies the cycle\item \textbf{"Was this the wrong moment?"}\item Three acts = the three assignment questions\end{itemize}
\smallskip{\color{navy}\textbf{\#}}\ inflation 11.6\% · gap −33 / −47pp\par
\vfill\textcolor{grey}{\rule{\linewidth}{0.3pt}}\par{\bfseries\color{navy}→}\ "First, the setting: the timeline."
\end{card}
\hfill
\begin{card}{navy}{\makebox[0pt][l]{5 \quad Timeline}\hfill{\normalsize\mdseries Zheng · 1:00 · from 4:30}}
{\large\itshape "The path was announced in March 2020."}\par\smallskip
\begin{itemize}[label=\textcolor{navy}{\textbullet}]\item Feb 2022: new framework\item May 2022: \textbf{1\%} · May 2023: \textbf{2\%} + O-SII cuts\item Both bind \textbf{31 May 2024}; held at 2\% ever since\item Point to the \textbf{gold band}: both decisions inside ECB hiking\end{itemize}
\smallskip{\color{navy}\textbf{\#}}\ last confirmed: Sep 2026, still 2\%\par
\vfill\textcolor{grey}{\rule{\linewidth}{0.3pt}}\par{\bfseries\color{navy}→}\ "What did the economy look like?"
\end{card}


\vspace{0.35cm}

\begin{card}{navy}{\makebox[0pt][l]{6 \quad Macro}\hfill{\normalsize\mdseries Zheng · 1:05 · from 5:30}}
{\large\itshape "The recovery was complete."}\par\smallskip
\begin{itemize}[label=\textcolor{navy}{\textbullet}]\item GDP above pre-COVID by \textbf{2021Q3}, ≈ +3\% by end-2021\item Then the storm: DNB quote, "outlook has worsened"\item Arguments \textbf{both ways}: build now vs wait\item Debt-sustainability warning\end{itemize}
\smallskip{\color{navy}\textbf{\#}}\ HICP 11.6\% (2022) · ECB +50bp, 21 Jul 2022\par
\vfill\textcolor{grey}{\rule{\linewidth}{0.3pt}}\par{\bfseries\color{navy}→}\ "In the financial system, the picture was split."
\end{card}
\hfill
\begin{card}{navy}{\makebox[0pt][l]{7 \quad Credit \& housing}\hfill{\normalsize\mdseries Zheng · 1:15 · from 6:35}}
{\large\itshape "Housing was hot, credit was cold."}\par\smallskip
\begin{itemize}[label=\textcolor{navy}{\textbullet}]\item House prices: DNB itself said "\textbf{overheated}"\item …then fell about 9\% real within a year\item Household debt high, but \textbf{falling}\item \textcolor{gold!80!black}{\textbf{DEF}} debt-service ratio, \textbf{lowest since 2005}\item Boom in house prices, \textbf{not} in credit\end{itemize}
\smallskip{\color{navy}\textbf{\#}}\ +19\% y/y (22Q1) · 16\% above 2007 peak · 107\% vs 58\% of GDP\par
\vfill\textcolor{grey}{\rule{\linewidth}{0.3pt}}\par{\bfseries\color{navy}→}\ "Who would carry the losses? The banks."
\end{card}

\newpage

\begin{card}{navy}{\makebox[0pt][l]{8 \quad Banks}\hfill{\normalsize\mdseries Zheng · 1:05 · from 7:50}}
{\large\itshape "The banks were in good shape."}\par\smallskip
\begin{itemize}[label=\textcolor{navy}{\textbullet}]\item \textcolor{gold!80!black}{\textbf{DEF}} CET1 = core equity / RWA\item Stress test: −3.8pp, still above minimum\item Rising rates \textbf{lifted profits} (remember this, it's the twist)\item Summary: \textbf{against} vs \textbf{for} raising now\end{itemize}
\smallskip{\color{navy}\textbf{\#}}\ CET1 17.7\% → 16.3\% · NII +6.7\%\par
\vfill\textcolor{grey}{\rule{\linewidth}{0.3pt}}\par{\bfseries\color{navy}→}\ \textbf{\textcolor{nlred}{\textbf{HAND-OVER:}}} "That was the setting. Diego will show what DNB looked at, and what it was insuring against."
\end{card}
\hfill
\begin{card}{nlred}{\makebox[0pt][l]{9 \quad Basel rule}\hfill{\normalsize\mdseries Diego · 1:15 · from 8:55}}
{\large\itshape "The rule DNB chose \emph{not} to follow."}\par\smallskip
\begin{itemize}[label=\textcolor{nlred}{\textbullet}]\item \textcolor{gold!80!black}{\textbf{DEF}} credit-to-GDP gap = distance from long-run trend (HP filter)\item Rule: < 2pp → 0\% · > 10pp → max\item \textbf{Missed} the GFC · \textbf{false alarm} 2012 · \textbf{rewritten} 2016\item Why: structural deepening, 47\% → 350\% of GDP; the filter calls it "trend"\end{itemize}
\smallskip{\color{nlred}\textbf{\#}}\ −13pp (2007) · +17pp (2012) · −0.6 vs +27pp (2016) · −33 / −47pp at the decisions\par
\vfill\textcolor{grey}{\rule{\linewidth}{0.3pt}}\par{\bfseries\color{nlred}→}\ "So DNB used a different compass."
\end{card}


\vspace{0.35cm}

\begin{card}{nlred}{\makebox[0pt][l]{10 \quad Framework}\hfill{\normalsize\mdseries Diego · 1:15 · from 10:10}}
{\large\itshape "DNB's 2022 framework has four phases."}\par\smallskip
\begin{itemize}[label=\textcolor{nlred}{\textbullet}]\item Recovery → \textbf{normal = 2\%} → elevated > 2\% → release\item \textcolor{gold!80!black}{\textbf{DEF}} positive neutral = filled when risk is "normal"\item Calibrated on \textbf{history}\item \textcolor{gold!80!black}{\textbf{DEF}} guided discretion = dashboard informs, doesn't dictate\item 17 jurisdictions; SE, UK, PL also 2\%\end{itemize}
\smallskip{\color{nlred}\textbf{\#}}\ €12bn losses → €6bn buffer → €150bn of lending\par
\vfill\textcolor{grey}{\rule{\linewidth}{0.3pt}}\par{\bfseries\color{nlred}→}\ "What reasons did DNB give?"
\end{card}
\hfill
\begin{card}{nlred}{\makebox[0pt][l]{11 \quad DNB's reasons}\hfill{\normalsize\mdseries Diego · 1:00 · from 11:25}}
{\large\itshape "What reasons did DNB give?"}\par\smallskip
\begin{itemize}[label=\textcolor{nlred}{\textbullet}]\item 2022: "\textbf{restore the buffers}", recovery, risks "normal to elevated"\item 2023: gap shows "\textbf{no signs of excessive credit growth}"\item Cited: risk appetite · falling real estate · debt · robust banks\item Never credit growth: \textbf{risk picture + strong banks + the promise}\end{itemize}
\vfill\textcolor{grey}{\rule{\linewidth}{0.3pt}}\par{\bfseries\color{nlred}→}\ "So which risks was it meant to address?"
\end{card}

\newpage

\begin{card}{nlred}{\makebox[0pt][l]{12 \quad Risks → tools}\hfill{\normalsize\mdseries Diego · 1:15 · from 12:25}}
{\large\itshape "Each risk was matched to a tool."}\par\smallskip
\begin{itemize}[label=\textcolor{nlred}{\textbullet}]\item \textbf{Unforeseeable shocks} → CCyB (war, energy, tightening, SVB/CS)\item \textbf{Housing} → LTV/LTI limits, tax, risk-weight floor\item \textbf{Big-bank size} → O-SII buffer\item \textcolor{gold!80!black}{\textbf{DEF}} Tinbergen = one instrument per target\item DNB: "sensitivity … to external events"\end{itemize}
\vfill\textcolor{grey}{\rule{\linewidth}{0.3pt}}\par{\bfseries\color{nlred}→}\ "Now we can test our two suspects."
\end{card}
\hfill
\begin{card}{nlred}{\makebox[0pt][l]{13 \quad Verdict ★}\hfill{\normalsize\mdseries Diego · 1:30 · from 13:40}}
{\large\itshape "Overheating or normalising? Row by row."}\par\smallskip
\begin{itemize}[label=\textcolor{nlred}{\textbullet}]\item Indicators \textbf{falling}, buffer up ✔\item +1pp a year, stop at 2 ✔\item 2023 prices falling: \textbf{raised anyway} ✔\item 2024–26 prices up \textasciitilde{}10\%: \textbf{held at 2} ✔\item Structural buffers \textbf{cut twice} ✔\item \emph{(pause)} "Five for five."\item Counter-evidence: housing +19\%, but DNB used other tools\end{itemize}
\vfill\textcolor{grey}{\rule{\linewidth}{0.3pt}}\par{\bfseries\color{nlred}→}\ "That last row is the key to the whole story."
\end{card}


\vspace{0.35cm}

\begin{card}{nlred}{\makebox[0pt][l]{14 \quad Swap}\hfill{\normalsize\mdseries Diego · 1:15 · from 15:10}}
{\large\itshape "What that last row means in euros."}\par\smallskip
\begin{itemize}[label=\textcolor{nlred}{\textbullet}]\item −6.0 (2020) · +3.3 · +3.4 · −2.7 (2024)\item Net vs pre-COVID ≈ \textbf{−2.0}, roughly unchanged\item DNB: "more or less capital-\textbf{neutral}" · Knot: "\textbf{limited} increase"\item Caveats: different bases; vs 2021 = an increase\item Capital \textbf{moved}, not added (the UK did the same)\end{itemize}
\vfill\textcolor{grey}{\rule{\linewidth}{0.3pt}}\par{\bfseries\color{nlred}→}\ "But a swap still has a cost. Was it the wrong moment?"
\end{card}
\hfill
\begin{card}{nlred}{\makebox[0pt][l]{15 \quad Twist ★}\hfill{\normalsize\mdseries Diego · 1:15 · from 16:25}}
{\large\itshape "The storm made it the \textbf{cheapest} moment."}\par\smallskip
\begin{itemize}[label=\textcolor{nlred}{\textbullet}]\item \textcolor{gold!80!black}{\textbf{DEF}} headroom = capital above requirements\item 2\% buffer ≤ \textbf{a fifth} of headroom (upper bound)\item Rates ↑ → profits ↑ → buffer from \textbf{retained earnings}\item ECB: gradual + profitable "limits costs"; during tightening "mitigates costs"\item \textbf{"Not the wrong moment. Arguably the right one."}\end{itemize}
\smallskip{\color{nlred}\textbf{\#}}\ €36bn headroom · €6.7bn = 19\% · ING, Rabo > €10bn each\par
\vfill\textcolor{grey}{\rule{\linewidth}{0.3pt}}\par{\bfseries\color{nlred}→}\ "Cheap for banks, but did borrowers pay?"
\end{card}

\newpage

\begin{card}{nlred}{\makebox[0pt][l]{16 \quad Borrowers}\hfill{\normalsize\mdseries Diego · 1:15 · from 17:40}}
{\large\itshape "Did borrowers pay?"}\par\smallskip
\begin{itemize}[label=\textcolor{nlred}{\textbullet}]\item \textcolor{gold!80!black}{\textbf{DEF}} event study, NL vs non-changers, around 3 dates\item \textcolor{gold!80!black}{\textbf{DEF}} placebo, each control as "fake treated" = grey band\item Rates \textbf{inside the band} · household credit grew \textbf{faster}\item Honest: NL lagged before (catch-up)\item Claim: "\textbf{no sign of contraction}", not causal\end{itemize}
\smallskip{\color{nlred}\textbf{\#}}\ credit 2022–26: +16.6\% vs +7.9\% (before: +4.3\% vs +11.3\%)\par
\vfill\textcolor{grey}{\rule{\linewidth}{0.3pt}}\par{\bfseries\color{nlred}→}\ \textbf{\textcolor{nlred}{\textbf{HAND-OVER:}}} "Cheap, and credit kept flowing. But a buffer reaches beyond Dutch banks. Zheng."
\end{card}
\hfill
\begin{card}{navy}{\makebox[0pt][l]{17 \quad Monetary policy}\hfill{\normalsize\mdseries Zheng · 1:05 · from 18:55}}
{\large\itshape "First wider implication: monetary policy."}\par\smallskip
\begin{itemize}[label=\textcolor{navy}{\textbullet}]\item Worry: \textbf{double squeeze} (ECB + DNB)\item Rates didn't diverge; profits rose\item ECB: early buffers let monetary policy "focus on price stability"\item \textcolor{gold!80!black}{\textbf{DEF}} macroprudential = stability of the whole system\item One ECB rate, \textbf{national} buffer = the Dutch lever + insurance\end{itemize}
\vfill\textcolor{grey}{\rule{\linewidth}{0.3pt}}\par{\bfseries\color{navy}→}\ "Second: borders."
\end{card}


\vspace{0.35cm}

\begin{card}{navy}{\makebox[0pt][l]{18 \quad Across borders}\hfill{\normalsize\mdseries Zheng · 1:00 · from 20:00}}
{\large\itshape "Second, borders."}\par\smallskip
\begin{itemize}[label=\textcolor{navy}{\textbullet}]\item \textcolor{gold!80!black}{\textbf{DEF}} reciprocity = foreign banks hold the Dutch 2\% too\item Limits leakage; non-bank leakage remains\item Europe moving the same way\item NL = highest neutral rate in the banking union (2023)\end{itemize}
\smallskip{\color{navy}\textbf{\#}}\ DE 0.75 · FR 1 · IE 1.5 · BE 1 · PT 0.75 · ES 1 (\%)\par
\vfill\textcolor{grey}{\rule{\linewidth}{0.3pt}}\par{\bfseries\color{navy}→}\ "Third: what the Netherlands gained."
\end{card}
\hfill
\begin{card}{navy}{\makebox[0pt][l]{19 \quad Resilience}\hfill{\normalsize\mdseries Zheng · 1:00 · from 21:00}}
{\large\itshape "What the Netherlands gained."}\par\smallskip
\begin{itemize}[label=\textcolor{navy}{\textbullet}]\item Releasable by design: \textbf{€0 → €6.7bn}, ≈ half of past peak losses\item Asymmetry: cheap to hold, powerful to release\item Floor expires \textbf{30 Nov 2026}: the CCyB matters more\end{itemize}
\smallskip{\color{navy}\textbf{\#}}\ up to €150bn lending · −0.1\% vs up to +10\%\par
\vfill\textcolor{grey}{\rule{\linewidth}{0.3pt}}\par{\bfseries\color{navy}→}\ "Fourth: it's not a clean win."
\end{card}

\newpage

\begin{card}{navy}{\makebox[0pt][l]{20 \quad Critiques}\hfill{\normalsize\mdseries Zheng · 1:00 · from 22:00}}
{\large\itshape "This is not a clean win."}\par\smallskip
\begin{itemize}[label=\textcolor{navy}{\textbullet}]\item \textbf{Credibility:} discretion must be explained, every quarter\item \textbf{Calibration:} why 2\%? ECB says 1.1–1.8\%\item \textbf{Untested:} lend or hoard?\item \textbf{Baseline:} vs 2021 it \emph{was} an increase\item \textbf{Profits:} excess-profit levy competed for the same money\end{itemize}
\vfill\textcolor{grey}{\rule{\linewidth}{0.3pt}}\par{\bfseries\color{navy}→}\ "Let me answer the three questions."
\end{card}
\hfill
\begin{card}{navy}{\makebox[0pt][l]{21 \quad Conclusion}\hfill{\normalsize\mdseries Zheng · 1:30 · from 23:00}}
{\large\itshape "Three answers."}\par\smallskip
\begin{itemize}[label=\textcolor{navy}{\textbullet}]\item \textbf{Environment:} recovery in a storm · credit cold · strong banks\item \textbf{Criteria:} 2\% normal level · past losses · any shock\item \textbf{Implications:} cheap · swap · no contraction · €6.7bn\item Not \emph{how much} capital, but \emph{which kind}\item \emph{(pause)} "…nothing to release. \textbf{Next time, it will.}"\item "Fill a buffer when nobody thinks you need it, even in a storm."\end{itemize}
\vfill\textcolor{grey}{\rule{\linewidth}{0.3pt}}\par{\bfseries\color{navy}→}\ "Thank you; we look forward to your questions."
\end{card}


\vspace{0.35cm}

\begin{card}{grey}{\makebox[0pt][l]{22 \quad Q\&A backup}\hfill{\normalsize\mdseries both · open · from 24:30}}
{\large\itshape Methods slide stays up during questions.}\par\smallskip
\begin{itemize}[label=\textcolor{grey}{\textbullet}]\item One speaker per question, about 45 seconds\item Default: Zheng answers policy, Diego answers econometrics\item Unknown? "We haven't measured that; our evidence is…"\item Always "\textbf{no sign of contraction}", never "no effect"\item Concede housing and the 2021 baseline first, then explain\end{itemize}
\end{card}
\hfill
\end{document}
